\documentclass[border=0pt]{standalone}
\usepackage[T1]{fontenc}
\usepackage{palatino}
\usepackage{graphicx}
\usepackage{tikz}
\usepackage{xcolor}

\begin{document}
\begin{tikzpicture}

% Black background
\fill[black] (0,0) rectangle (14.437348000000002in, 10.417in);

% Background image
\begin{scope}
\clip (0,0) rectangle (14.437348000000002in, 10.417in);
\node[inner sep=0pt, outer sep=0pt] at (3.0003638333333322in, 5.2085in) {
  \includegraphics[height=10.417in]{../figures/art-consciousness-ultimate-upscale.png}
};
\end{scope}

% ---- FRONT COVER (right side) ----

% Dark gradient at top — covers title + subtitle, fades before neural eye
\fill[top color=black!90, bottom color=black!0, opacity=0.7]
  (7.752348in, 7.08356in)
  rectangle (14.437348000000002in, 10.417in);

% Gradient at bottom for author
\fill[bottom color=black!80, top color=black!0, opacity=0.65]
  (7.752348in, 0in)
  rectangle (14.437348000000002in, 1.9084999999999999in);

% Front title
\node[anchor=north, text=white, font=\fontsize{36}{42}\selectfont\bfseries,
      text width=5.0in, align=center]
  at (10.752348in, 9.358500000000001in)
  {Die Simulation\\[0.15in] namens Ich};

% Front subtitle (within gradient zone)
\node[anchor=north, text=white!90, font=\fontsize{14}{18}\selectfont,
      text width=4.8in, align=center]
  at (10.752348in, 8.008500000000002in)
  {Die Architektur von Bewusstsein,\\Berechnung und Kosmos};

% Front author
\node[anchor=south, text=white!95, font=\fontsize{18}{22}\selectfont]
  at (10.752348in, 1.2085in)
  {Matthias Gruber};

% ---- SPINE ----

% Dark overlay on spine + hinge zones for readability
\fill[black, opacity=0.5]
  (6.882in, 0in)
  rectangle (7.5553479999999995in, 10.417in);

% Spine text (rotated)
\node[rotate=270, text=white, font=\fontsize{10}{12}\selectfont\bfseries,
      anchor=center]
  at (7.218674in, 5.7085in)
  {Die Simulation namens Ich};

\node[rotate=270, text=white!90, font=\fontsize{8}{10}\selectfont,
      anchor=center]
  at (7.218674in, 3.7085in)
  {Matthias Gruber};

% ---- BACK COVER (left side) ----

% Dark overlay on back cover for text
\fill[black, opacity=0.55]
  (0.6849999999999999in, 0.7084999999999999in)
  rectangle (6.685in, 9.7085in);

% Back cover blurb
\node[anchor=north, text=white!95, font=\fontsize{11}{15}\selectfont,
      text width=4.6in, align=left]
  at (3.685in, 8.7085in)
  {Das Ich ist eine Simulation.

Keine Metapher. Kein philosophisches Gedankenexperiment. Eine buchstäblich laufende Simulation ---
vom Gehirn erzeugt, hunderte Male pro Sekunde aktualisiert, von innen erlebt.

Dieses Buch entwickelt eine einheitliche Theorie des Bewusstseins, aufgebaut auf vier Modellen,
die das Gehirn unterhält: zwei der Welt (eines gelernt, eines simuliert) und zwei des Selbst
(eines gelernt, eines simuliert). Die Theorie löst das \quotedblbase{}Schwierige Problem\textquotedblleft{} auf, erklärt die
Phänomenologie psychedelischer Zustände, macht neun falsifizierbare Vorhersagen und liefert
einen konkreten Bauplan für eine bewusste Maschine.

Wenn die Theorie stimmt, war Bewusstsein nie ein Rätsel. Es war nur falsch etikettiert.};

% ISBN barcode area (white box — KDP overlays free ISBN barcode here)
\fill[white] (4.484999999999999in, 1.3185in)
  rectangle (6.484999999999999in, 2.5185in);
\node[inner sep=0pt] at (5.484999999999999in, 1.9184999999999999in)
  {\includegraphics[width=1.8in]{isbn-barcode-hc-de.png}};

\end{tikzpicture}
\end{document}
